\scanpage{94}

\textcircled{5.} \uline{Classes remarquables d'opérateurs dans l'espace de Hilbert}

Nous commençons par donner un exposé de la théorie de Schatten, \deleted{\unsure{qui est}} sous une forme simplifiée rendue possible par la \deleted{\unsure{lemme}}\inserted{méthode du \unsure{minimum|minimax|convexité}} de Weyl, et en y apportant quelques compléments.

\margin{?}

\deleted{\unsure{Déf}} Rappelons la

\uline{Définition 1} \deleted{Soit} Une norme $N$ sur $R^n$ est dite \uline{norme de Schatten} si elle est symétrique, satisfait à \deleted{\unsure{à}}

\[
N(x_1,\ldots x_n)=N(|x_1|,\ldots |x_n|)\tag*{(1)}
\]

\deleted{et si elle} si elle est croissante \deleted{\unsure{sur l'ensem}} sur l'ens. des $x_i\geq 0$, et \deleted{\unsure{est satisf|elle satisfait}} est normalisée par $N(1,0,\ldots 0)=1$

Une norme $N$ sur $R^{(\underline N)}$ \deleted{\unsure{$\underline N$ ensemble}} (\unsure{espace|anneau} des suites infinies \deleted{\unsure{finies}} nulles sauf un nb fini de termes) est appellée une norme de Schatten, si \inserted{pour tout entier $n>0$,} elle induit \deleted{sur} \deleted{de} sur $R^n$ une norme de Schatten

($R^n$ sera ici identifié \deleted{à l'espace des suites}\inserted{au sous-espace} de $R^{(\underline N)}$ formé des suites nulles à partir du rang $n+1$). Notons qu'une norme de Schatten sur $R^n$ induit sur $R^m$ ($m\leq n$) une norme de Schatten. \deleted{\unsure{Tout cette théorie résulte que}}\inserted{Exemples de normes de} Schatten\,: \deleted{\unsure{Supposons}}

\[
\deleted{N_p(\{x_i\})=\left(\sum |x_i|^p\right)^{1/p}\quad\text{où }1\leq p}
\]

les normes $N_p((x_i))=\lVert(x_i)\rVert_p$, définies par \deleted{$1\leq p\leq\infty$ \unsure{par}}

\[
\lVert(x_i)\rVert_p=\left(\sum |x_i|^p\right)^{1/p}\qquad (1\leq p<+\infty)\tag*{(2)}
\]

\[
\lVert(x_i)\rVert_\infty=\underset{i}{\operatorname{Sup}}\,|x_i|\tag*{\deleted{\unsure{4}} (3)}
\]

Si $N_0$ est une norme de Schatten \deleted{alors} sur $R^n$, il existe toujours une norme de Schatten $N$ sur $R^{(\underline N)}$ \deleted{\unsure{qui}} induise $N_0$\deleted{, \unsure{p. ex.,} \illegible{1}} si, pour \deleted{\unsure{toute|tout}} $x=(x_i)\in\underline{c}_0$, on désigne par $(\unsure{\rho|f|R}_i(x))$ la suite des $|x_i|$ \deleted{\unsure{rangés}}\inserted{mais} rangés par ordre de grandeur décroissant, on peut \deleted{\unsure{p.ex.}} poser $N(x)=N_0(\rho_1(x),\ldots\rho_n(x))$. Il est trivial

\margin{\unsure{?}}
% archive sheet number: 95
